\section{C 语言要点（嵌入式视角）}

裸机开发里，C 语言的一些特性会被高频使用，理解它们有助于读懂 CMSIS 和
标准外设库源码。本节挑重点讲。

\subsection{固定宽度整数：stdint.h}

嵌入式里寄存器宽度是硬件定死的，必须用固定宽度类型而不是 \texttt{int}：

\begin{lstlisting}[style=code]
#include <stdint.h>

uint32_t value;   // 一定是 32 位无符号
uint16_t reg16;   // 16 位
uint8_t  flag;    // 8 位
int32_t  signed32;
\end{lstlisting}

\texttt{int} 的宽度随平台/编译器而变，在寄存器操作里用 \texttt{int} 是隐患。

\subsection{volatile：防止编译器优化掉内存访问}

外设寄存器是\textbf{内存映射}的，硬件会异步修改它们，必须用
\texttt{volatile} 告诉编译器「每次都要真正读写内存，不能缓存、不能省略」：

\begin{lstlisting}[style=code]
volatile uint32_t *reg = (volatile uint32_t *)0x40023800;
*reg = 0x01;       // 必须真正写到该地址
while (*reg & 0x02) { /* 必须每次重新读 */ }
\end{lstlisting}

漏掉 \texttt{volatile} 的典型后果是：编译器把循环里的读取优化掉，导致忙等死循环。

\subsection{CMSIS 的寄存器访问模型}

CMSIS 不直接写裸指针，而是用一个\textbf{结构体覆盖}外设地址空间，
并用 \texttt{\_\_IO}/\texttt{\_\_IM}/\texttt{\_\_OM} 标注读写属性：

\begin{lstlisting}[style=code]
/* __IO = volatile 可读可写，__IM = 只读，__OM = 只写 */
typedef struct {
  __IO uint32_t MODER;    /* 模式寄存器 */
  __IO uint32_t OTYPER;
  __IO uint32_t OSPEEDR;
  __IO uint32_t PUPDR;
  __IO uint32_t IDR;
  __IO uint32_t ODR;
} GPIO_TypeDef;

#define GPIOD ((GPIO_TypeDef *)0x40020C00)  /* 外设基地址 */
GPIOD->ODR |= (1U << 12);                    /* 直接操作寄存器 */
\end{lstlisting}

这种写法类型安全、可读性好，是 STM32 库（HAL/StdPeriph/CMSIS）的通用套路。

\subsection{位操作}

\begin{lstlisting}[style=code]
#define BIT(n) (1UL << (n))

uint32_t r = 0;

r |=  BIT(12);      // 置位：第 12 位 = 1
r &= ~BIT(12);      // 清位：第 12 位 = 0
r ^=  BIT(12);      // 翻转
if (r & BIT(12)) {} // 测试

// 写多位（例如写 MODE 字段，占 2 位）
r = (r & ~(0x3UL << 10)) | (value << 10);   // 先清空字段，再写入
\end{lstlisting}

「先清空、再写入」是操作寄存器字段的标准范式，库里的
\texttt{GPIO\_Init} 等函数内部正是这么写的。

\subsection{freestanding 与 -nostdlib}

标准 C 区分「宿主环境（hosted）」和「独立环境（freestanding）」。
裸机没有操作系统，属于 freestanding：

\begin{itemize}
  \item 编译选项加了 \texttt{-ffreestanding}，告诉编译器不假设有 \texttt{libc}；
  \item 链接选项加了 \texttt{-nostdlib -nostartfiles}，不链接标准库和启动文件；
  \item 因此 \texttt{printf}/\texttt{malloc}/\texttt{memcpy} 等\textbf{默认不可用}。
        需要时要么自己实现，要么链接 newlib 并配置 \texttt{syscalls}。
\end{itemize}

这也是本工程在链接脚本/启动文件里去掉 \texttt{libc.a} 和
\texttt{\_\_libc\_init\_array} 的根本原因。

\subsection{weak 弱符号与中断处理函数}

启动文件里每个中断向量都指向一个\textbf{弱符号}的默认处理函数：

\begin{lstlisting}[style=generic]
.weak  WWDG_IRQHandler
.thumb_set WWDG_IRQHandler, Default_Handler
\end{lstlisting}

\texttt{.weak} 表示「如果你在别处定义了同名函数，就用你的；否则用默认的」。
所以用户在 \texttt{main.c} 里写一个 \texttt{void TIM2\_IRQHandler(void)} 就能
覆盖默认实现、接管中断，无需改启动文件。

\subsection{const 与 flash 布局}

\begin{lstlisting}[style=code]
const char msg[] = "hello";   // 放到 .rodata（Flash），不占 RAM
\end{lstlisting}

裸机 RAM 很小，常量应尽量用 \texttt{const} 让链接器把它们放进 Flash 的
只读段，而不是可写的 \texttt{.data}。
